% =========================================================
% Hoja de vida de EJEMPLO (candidato ficticio).
% Cópiala a cv/plantillas/hoja_de_vida_<rol>.tex y reemplaza todo con tus datos reales.
% cv/plantillas/ está en .gitignore: tus hojas de vida y tu foto nunca se suben a git.
% Compila con: pdflatex -interaction=nonstopmode -halt-on-error hoja_de_vida_<rol>.tex
% Debe quedar en exactamente 1 página.
% =========================================================

\documentclass[10pt,a4paper]{article}

% ---------- Interruptores ----------
% \fototrue necesita Foto_perfil.jpeg junto a este archivo.
\newif\iffoto       \fotofalse
% \anexostrue agrega al final los PDF de soporte listados en la sección Anexos.
\newif\ifanexos     \anexosfalse

\usepackage{pdfpages}

% ---------- Paquetes ----------
\usepackage[top=0.8cm,bottom=0.6cm,left=1.5cm,right=1.5cm]{geometry}
\usepackage[T1]{fontenc}
\usepackage[utf8]{inputenc}
\IfFileExists{newtxtext.sty}{\usepackage{newtxtext}}{\usepackage{mathptmx}}
\usepackage{graphicx,tikz,xcolor,enumitem,titlesec}
\usepackage[hidelinks]{hyperref}

% Texto extraíble correctamente por los ATS (ligaduras y tildes)
\input{glyphtounicode}
\pdfgentounicode=1

% Sin cortes de palabra con guion
\hyphenpenalty=10000
\exhyphenpenalty=10000
\sloppy

% El dashboard (Buscador) lee el rol del final de pdftitle y las palabras clave de pdfkeywords.
\hypersetup{
  pdftitle={Nombre Apellido - Hoja de Vida - Analista de Datos Junior},
  pdfauthor={Nombre Apellido},
  pdfsubject={Analista de Datos Junior: SQL, Power BI, Excel avanzado, indicadores de gestión y calidad de datos},
  pdfkeywords={Analista de datos, análisis de datos, SQL, PostgreSQL, Power BI, Excel avanzado, Power Query, indicadores de gestión, KPI, tableros, reportes, calidad de datos, Python, pandas}
}

% ---------- Estilo ----------
\pagestyle{empty}
\setlength{\parindent}{0pt}
\definecolor{acento}{RGB}{31,56,100}
\definecolor{linea}{RGB}{150,150,150}

\titleformat{\section}{\large\bfseries\color{acento}}{}{0pt}{\MakeUppercase}[{\color{linea}\titlerule[0.5pt]}]
\titlespacing*{\section}{0pt}{4pt}{2pt}

\setlist[itemize]{leftmargin=1.1em,label=\textbullet,topsep=2pt,itemsep=0.5pt,parsep=0pt}

% \cargo{Cargo}{Empresa}{Ciudad}{Fechas}
\newcommand{\cargo}[4]{%
  \par\noindent\textbf{#1} \textbar\ #4\par
  \noindent\textit{#2} · #3\par}

\begin{document}
\small\linespread{0.97}\selectfont

% ---------- Encabezado ----------
\noindent
\begin{minipage}[c]{\iffoto 0.83\textwidth\else\textwidth\fi}
  {\fontsize{20}{22}\selectfont\bfseries Nombre Apellido}\par\vspace{3pt}
  {\normalsize\bfseries\color{acento} Analista de Datos \textbar\ SQL, Power BI, Excel e indicadores de gestión}\par
  \vspace{4pt}
  Ciudad, Departamento, País \textbar\ +57 300 000 0000 \textbar\ Disponibilidad inmediata\par
  \href{mailto:nombre@example.com}{nombre@example.com} \textbar\ \href{https://github.com/usuario}{github.com/usuario}\par
  \href{https://www.linkedin.com/in/usuario}{linkedin.com/in/usuario}\par
\end{minipage}%
\iffoto
\hfill
\begin{minipage}[c]{2.3cm}
  \begin{tikzpicture}
    \clip (0,0) circle (1.12cm);
    \node[inner sep=0pt] at (0,-0.12cm){\includegraphics[width=2.55cm]{Foto_perfil.jpeg}};
  \end{tikzpicture}
\end{minipage}
\fi
\par

% ---------- Perfil ----------
\section{Perfil profesional}
Profesional en Ingeniería de Sistemas con dos años en análisis de datos, reportería y construcción de indicadores de gestión. Escribo consultas SQL sobre PostgreSQL y construyo tableros en Power BI y Excel avanzado con Power Query. Valido y depuro información de varias fuentes para entregar datos consistentes. Este párrafo es un ejemplo: describe aquí solo lo que puedas demostrar.

% ---------- Experiencia ----------
\section{Experiencia profesional}

\cargo{Analista de Datos}{Empresa Ejemplo S.A.S.}{Ciudad}{01/2024 -- Actualidad}
\begin{itemize}
  \item Construyo y mantengo los tableros e indicadores de gestión (KPI) de ventas que usa la gerencia, sobre consultas SQL en PostgreSQL.
  \item Automaticé la consolidación semanal de reportes en Excel con Power Query, lo que eliminó el trabajo manual de copiar datos.
  \item Valido y depuro datos de tres sistemas de origen antes de publicarlos. % CONFIRMAR: número real de sistemas de origen.
\end{itemize}\vspace{3pt}

\cargo{Auxiliar de Información}{Comercializadora Ficticia Ltda.}{Ciudad}{02/2022 -- 12/2023}
\begin{itemize}
  \item Digitalicé y mantuve actualizadas las bases de datos de clientes e inventario.
  \item Elaboré reportes mensuales en Excel y Power BI para el área comercial.
  \item Brindé soporte a los usuarios en Microsoft Excel y Word.
\end{itemize}

% ---------- Proyectos ----------
\section{Proyectos destacados}
\begin{itemize}\raggedright
  \item \textbf{Tablero de ventas de ejemplo} (PostgreSQL, SQL, Power BI). Modelo de datos de ventas e inventario con un tablero de indicadores por sucursal. \href{https://github.com/usuario/proyecto-ejemplo}{github.com/usuario/proyecto-ejemplo}
  \item \textbf{Proyecto de grado: análisis de datos de producción} (Python, pandas). Limpieza de registros de producción y tablero de indicadores de la línea.
\end{itemize}

% ---------- Habilidades ----------
\section{Habilidades técnicas}
\textbf{Reportería y BI:} Power BI, Excel avanzado, Power Query; tableros e indicadores de gestión (KPI).\par
\textbf{SQL y bases de datos:} consultas SQL, PostgreSQL, MySQL, modelado de datos relacional.\par
\textbf{Calidad de datos:} validación y depuración de información de varias fuentes; limpieza con Python (pandas).\par
\textbf{Otras herramientas:} Git, GitHub, Linux.\par

% ---------- Educación ----------
\section{Educación}
\cargo{Ingeniero de Sistemas}{Universidad Ejemplo}{Ciudad}{2018 -- 2023}
Formación en bases de datos, estadística y programación.\par

% ---------- Certificaciones ----------
\section{Cursos, certificaciones e idiomas}
\begin{itemize}[itemsep=0pt]
  \item Curso de ejemplo en visualización de datos, Plataforma Ejemplo (01/2025)
  \item Curso de ejemplo en bases de datos con SQL, Institución Ejemplo (06/2024)
\end{itemize}
\vspace{2pt}\noindent\textbf{Idiomas:} español nativo \textbar\ inglés intermedio

% ---------- Anexos ----------
\ifanexos
\newpage
% Activa \anexostrue solo cuando estos PDF existan junto a la hoja de vida.
\includepdf[pages=-]{anexos/certificado_ejemplo.pdf}
\fi

\end{document}
